%-----------------------------------------------------------------
% SURAT PERNYATAAN
%
% Seluruh data diisi otomatis dari include/identitas.tex
%-----------------------------------------------------------------
\spernyataan % \spernyataan JANGAN DIHAPUS!

\noindent
Yang bertanda tangan di bawah ini,
\vspace{-0.1cm}

\begin{table}[H]
\renewcommand{\arraystretch}{0.75}
\begin{tabularx}{\textwidth}{@{}l p{2.9cm} c >{\raggedright\arraybackslash}X@{}}
	1.	& Nama   	& : & \fullname \\
		& NIM   	& : & \idnum \\
		& Program Studi	& : & \programstudi \\
		& Status 	& : & \statusmahasiswa \\[4pt]
	2.	& Nama   	& : & \firstsupervisor \\
		& NIP    	& : & \firstnip \\
		& Program Studi	& : & \programstudi \\
		& Status 	& : & Pembimbing I \\[4pt]
	3.	& Nama   	& : & \secondsupervisor \\
		& NIP    	& : & \secondnip \\
		& Program Studi	& : & \programstudi \\
		& Status 	& : & Pembimbing II
\end{tabularx}
\end{table}

\vspace{-0.2cm}
\noindent
Dengan ini menyatakan hasil penelitian Tesis yang berjudul
\textbf{``\MakeUppercase{\judul}''} tidak dipublikasikan secara
\textit{full-text} di sistem ETD (\textit{Electronic Theses and
Dissertations}) Universitas Syiah Kuala hingga batas waktu 5 tahun dari
tanggal kelulusan.

\vspace{0.3cm}
\noindent
Demikian surat pernyataan ini dibuat dengan sebenarnya untuk dapat
dipergunakan seperlunya.

\vspace{0.3cm}

{\centering
\city, \disahkantanggal\par
\vspace{0.15cm}
Yang membuat pernyataan,\par
}

%-------------------------------------------------------------
% Lebar kolom diukur lebih dulu pada blok BAWAH (Ketua Program
% Studi + Koordinator Tesis). Lebar itu lalu jadi acuan blok tiga
% kolom di atasnya: kolom Pembimbing I dan II ikut lebar mengikuti,
% sedangkan kolom mahasiswa dibiarkan setipis mungkin (namanya
% tanpa gelar) dan terdorong sampai ke tepi kanan. Jadi kalau
% blok bawah dilebarkan, blok atas ikut menyesuaikan.
%-------------------------------------------------------------
\ttdgutter=0.6cm
\ttdkolomA=0pt \ttdkolomB=0pt
\ttdukurA{Ketua Program Studi}\ttdukurA{\programstudi{} \fakultassingkat}%
\ttdukurA{\university,}\ttdukurA{\kaprodi}\ttdukurA{NIP. \kaprodinip}%
\ttdukurB{Koordinator Tesis,}\ttdukurB{\koordinatortesis}%
\ttdukurB{NIP. \koordinatortesisnip}%
\ttdfit
\newlength\ttlebarblokTiga
\setlength{\ttlebarblokTiga}{\dimexpr\ttdkolomA+\ttdkolomB+\ttdgutter\relax}

%-------------------------------------------------------------
% Tanda tangan mahasiswa dan kedua dosen pembimbing
%-------------------------------------------------------------
% Kolom mahasiswa dibiarkan pada lebar minimumnya (yang dikejar
% NIM), sedangkan kolom pembimbing I dan II memakai seluruh sisa
% ruang. Ruang blok tiga kolom diberi sedikit lebih longgar dari
% blok bawah supaya kedua kolom pembimbing benar-benar bisa
% melebar, tanpa keluar dari kertas.
\ttdgutter=0.35cm
\newlength\ttlebarAtas
\setlength{\ttlebarAtas}{\dimexpr\textwidth+45pt\relax}
\newlength\ttdgeserAtas
\ttdcolA=0pt \ttdcolB=0pt \ttdcolC=0pt
\ttdasarA{Pembimbing I,}\ttdasarA{NIP. \firstnip}\ttdnamaA{\firstsupervisor}%
\ttdasarB{Pembimbing II,}\ttdasarB{NIP. \secondnip}\ttdnamaB{\secondsupervisor}%
\ttdasarC{Mahasiswa,}\ttdasarC{NIM. \idnum}\ttdnamaC{\fullname}%
\ttdfitTigaPada{\ttlebarAtas}%
\setlength{\ttdgeserAtas}{\dimexpr(\ttlebarAtas-\textwidth)/2\relax}%
\ifdim\ttdgeserAtas>25pt \setlength{\ttdgeserAtas}{25pt}\fi
\par
\noindent\hspace*{\dimexpr-\ttdgeserAtas\relax}%
\parbox[t]{\ttdcolA}{\raggedright\setlength{\baselineskip}{14pt}%
  Pembimbing I,\par
  \vspace{1.1cm}%
  \ttdnamaW[\ttdcolA]{\firstsupervisor}\par
  \vspace{0.15cm}%
  NIP. \firstnip}%
\hspace*{\ttdgutter}%
\parbox[t]{\ttdcolB}{\raggedright\setlength{\baselineskip}{14pt}%
  Pembimbing II,\par
  \vspace{1.1cm}%
  \ttdnamaW[\ttdcolB]{\secondsupervisor}\par
  \vspace{0.15cm}%
  NIP. \secondnip}%
\hspace*{\ttdgutter}%
\parbox[t]{\ttdcolC}{\raggedright\setlength{\baselineskip}{14pt}%
  Mahasiswa,\par
  \vspace{1.1cm}%
  \ttdnamaW[\ttdcolC]{\fullname}\par
  \vspace{0.15cm}%
  NIM. \idnum}\par

\vspace{0.3cm}

%-------------------------------------------------------------
% Tanda tangan Ketua Program Studi dan Koordinator Tesis
%-------------------------------------------------------------
{\centering Mengetahui,\par}

\vspace{0.55cm}
\noindent\hspace*{\dimexpr-\ttdgeser\relax}%
\parbox[t]{\ttdkolomA}{\raggedright\setlength{\baselineskip}{14pt}%
  Ketua Program Studi\par
  \programstudi{} \fakultassingkat\par
  \university,}%
\hspace*{\ttdgutter}%
\parbox[t]{\ttdkolomB}{\raggedright\setlength{\baselineskip}{14pt}%
  Koordinator Tesis,}\par
\vspace{1.0cm}
\noindent\hspace*{\dimexpr-\ttdgeser\relax}%
\parbox[t]{\ttdkolomA}{\raggedright\setlength{\baselineskip}{14pt}%
  \ttdnama[\ttdkolomA]{\kaprodi}\par
  \vspace{0.15cm}%
  NIP. \kaprodinip}%
\hspace*{\ttdgutter}%
\parbox[t]{\ttdkolomB}{\raggedright\setlength{\baselineskip}{14pt}%
  \ttdnama[\ttdkolomB]{\koordinatortesis}\par
  \vspace{0.15cm}%
  NIP. \koordinatortesisnip}